% SHARD-ID: CA-06-SECTORS-PROJECTIONS-QUOTIENTS
% SHARD-TITLE: Sectors, Projections, Subspaces, and Quotients
% SHARD-SUMMARY: Makes precise the slogan that symmetry sectors can be represented by orthogonal projections, closed subspaces, or suitable Hilbert quotients.
% SHARD-SUMMARY: Proves the finite-group averaging projection for invariant sectors and records the checked C2 swap example in the Julia test suite.
% SHARD-SUMMARY: Separates invariant vectors, invariant observable algebras, and compressed sector observable algebras so the compiler cannot conflate them.
% SHARD-KEYWORDS: sector, projection, quotient, invariant subspace, averaging projector, observable algebra, compression

\section{Sectors, Projections, Subspaces, and Quotients}
\label{sec:sectors-projections-quotients}

\subsection{Sources and Dependencies}

This shard depends on CA-04 and CA-05 and fixes the first sector semantics for
the Hilbert-space compiler.  It uses CONVENTIONS.md (g).

% Source: references/text/PenneysUnitaryFusionCategories.md:219
%   orthogonal direct sums give isometries and mutually orthogonal projections.
% Source: references/text/PenneysUnitaryFusionCategories.md:235
%   an orthogonal projection satisfies p^2 = p = p^\dagger and splits by an isometry.
% Checked: test/runtests.jl, "finite symmetry sector projectors".

\subsection{The Triangle}

Let \(H\) be a Hilbert space.  The compiler treats a sector first as an
orthogonal projection
\[
  P=P^*=P^2\in\BH(H).
\]
It may then expose either of the equivalent Hilbert-space presentations
\[
  H_P:=\Ran(P)\subset H,\qquad
  H/\ker(P)=H/\Ran(1-P).
\]
The map
\[
  [v]\in H/\ker(P)\longmapsto Pv\in \Ran(P)
\]
is a well-defined unitary isomorphism when the quotient has the Hilbert quotient
norm.  Thus the safe slogan is not ``symmetry equals quotient'' but:
\[
  \text{projection}\quad\Longleftrightarrow\quad
  \text{closed subspace}\quad\Longleftrightarrow\quad
  \text{quotient by the named orthogonal complement}.
\]

\subsection{Finite-Group Invariants}

Let \(G\) be finite and \(U:G\to U(H)\) a unitary representation.  Define
\[
  P_G=\frac1{|G|}\sum_{g\in G}U(g).
\]
Then \(P_G\) is the orthogonal projection onto the invariant subspace
\[
  H^G=\{v\in H:U(g)v=v\ \forall g\in G\}.
\]

\emph{Proof.}  Idempotence follows by reindexing:
\[
  P_G^2=\frac1{|G|^2}\sum_{g,h}U(gh)
       =\frac1{|G|}\sum_k U(k)=P_G.
\]
Self-adjointness follows from unitarity and \(g\mapsto g^{-1}\):
\[
  P_G^*=\frac1{|G|}\sum_g U(g)^*
       =\frac1{|G|}\sum_g U(g^{-1})=P_G.
\]
For every \(k\), \(U(k)P_G=P_G\), so \(\Ran(P_G)\subset H^G\).  If \(v\in H^G\),
then \(P_Gv=v\), so \(H^G\subset\Ran(P_G)\). \(\square\)

\subsection{Coinvariant Quotient}

For the same finite action, let
\[
  R_G=\Span\{U(g)v-v:g\in G,\ v\in H\}.
\]
Then \(R_G\subseteq\ker(P_G)\) because \(P_GU(g)=P_G\).  Conversely, if
\(w\in\ker(P_G)\), then
\[
  w
  =
  \frac1{|G|}\sum_{g\in G}\bigl(w-U(g)w\bigr)\in R_G.
\]
Therefore \(R_G=\ker(P_G)=(H^G)^\perp\), and the coinvariant quotient
\[
  H_G:=H/R_G
\]
is canonically unitarily isomorphic to the invariant subspace \(H^G\).  For
compact non-finite groups the analogous statement requires a Haar-integral and
analytic hypotheses; this shard records only the finite checked case.

\subsection{Observable-Algebra Distinctions}

There are three different symmetry-related algebras that the compiler must not
identify silently:
\[
  \BH(H),\qquad
  \BH(H)^G:=\{A:U(g)AU(g)^*=A\ \forall g\},\qquad
  P_G\BH(H)P_G\cong \BH(H^G).
\]
The first is the full algebra before imposing symmetry.  The second is the
algebra of invariant observables on the original carrier.  The third is the
compressed algebra acting on the invariant sector.  A model may choose one of
these, but the choice is an observable policy in the sense of CA-04, not a
consequence of the word ``symmetry'' alone.

\subsection{Checked Example}

The current Julia check uses \(G=C_2\) acting on \(\mathbb C^2\) by the swap
matrix \(S\).  The averaging projection is
\[
  P=\frac12(I+S)
   =
  \begin{pmatrix}
    1/2 & 1/2\\
    1/2 & 1/2
  \end{pmatrix}.
\]
The test checks \(P^*=P=P^2\), checks \(I-P\) is the complementary orthogonal
projection, verifies \(P(I-P)=0\), and verifies that \((1,1)\) is fixed while
\((1,-1)\) is killed.  This is the first executable invariant for the sector
compiler.

\subsection{Lamport Dependency Skeleton}

{\small
\[
\setlength{\arraycolsep}{2pt}
\begin{array}{rcl}
D17 &:& \text{sector projection }P=P^*=P^2;\quad H_P=\Ran(P).\\
D18 &:& H/\ker(P)\xrightarrow{\ [v]\mapsto Pv\ }\Ran(P).\\
D19 &:& P_G=|G|^{-1}\sum_g U(g)\quad (G\text{ finite}).\\
D20 &:& R_G=\Span\{U(g)v-v\}.\\
L1 &:& P_G=P_G^*=P_G^2,\quad \Ran(P_G)=H^G.\\
L2 &:& R_G=\ker(P_G)=(H^G)^\perp,\quad H/R_G\cong H^G.\\
T11 &:& \text{Projection, subspace, and quotient are equivalent only after the relation subspace is named.}\\
T12 &:& \BH(H)^G\text{ and }P_G\BH(H)P_G\text{ are different observable policies.}
\end{array}
\]
}

\subsection{Boundary}

The finite averaging theorem is not a statement about spontaneous symmetry
breaking, gauge constraints, BRST cohomology, braid statistics, or superselection
sectors in a continuum CFT.  It is the first Hilbert-space sector rule needed by
the compiler.
